\section{Experimental chronology}
\label{sec:history}

Development began in June 2026 with the E-JEPA-TTC hypothesis: predictive latent
learning might improve event-based TTC estimation. The research then moved
through geometry and cross-time transport to temporal context and RGB--event
fusion. Several evaluation corrections changed which results could support a
claim.

\subsection{June--July: predictive learning and evaluation corrections}

The first experiments established indexed EvTTC access, voxel representations,
a TinyCNN, and geometric/event-rate baselines. JEPA pretraining with
latent-variance regularization led to multihorizon prediction, frozen-encoder
probes, and low-label comparisons. Token transformers, tubelets, deep token
supervision, navigation conditioning, and rollout probes tested other forms of
temporal information.

Centered derivatives gave the geometric baseline access to future observations.
We replaced them with past-only derivatives and reported this annotation-assisted
baseline on its supported queries, separate from models with different inputs or
coverage.

On 13 July, we reclassified the inspected CPLA-high results as reused-test
diagnostics (\texttt{26480260}). Their three fine-tuning seeds shared one SSL
initialization, so the spread reflects downstream optimization conditional on
that parent. We removed the runs from final-test claims and added checkpoint,
cache-version, nested-subset, and matched-scratch checks. No repartition of the
inspected sequences could recreate a blind holdout.

An eAP smoke-validation defect wrote success markers without running the matrix
(\texttt{1f9e2f99}). We replaced the faulty path with the matrix invocation and
added a regression test; the old markers were discarded as evidence.

\subsection{August: geometry, transport and fold-local parents}

August introduced physics-constrained FlowMimic, dense LevelDynamics,
object-centric representations, and matched Garl comparisons. The A-series
varied foreground geometry and decoder capacity, then added SAM supervision and
DINOv3 relational distillation. A4 used an RGB teacher; A5 learned local
cross-time transport on disjoint training rows. These stages no longer tested
the original label-free JEPA objective alone.

On 13 August, we found that several grouped runs inherited an A4 parent trained
on their held-out outer groups (\texttt{de3349fd}). Those summaries remain
parent-exposed diagnostics. The corrected study trained A4 inside each outer
fold, keeping the held-out group out of both representation learning and
downstream fitting.

V6 and V7 produced no promoted model. A post-hoc A5/C2F routing result from V7
became a prospective V8 hypothesis. V8 then tested temporal mechanisms and
nested routing after correcting ROI quantization, timestamp bins, and bounded
HDF5 reads. Unsupported predictions stayed in the evaluation as named failures.

\subsection{September--October: context and multimodal measurement}

X0/X1 and Stages 61--76 tested phase dynamics, raw timing, routing, risk, and
context under fixed optimization budgets. Nested producers excluded held-out
groups from training, and incomplete runs were not promoted to completed
evaluations.

SIMPLEX-T completed 24 head fits and 60,000 updates. Its H8/H16 comparisons
retrained heads while sharing upstream producers. Efficient-context work then
addressed TRAIN40 inputs and execution cost; the RGB port added native sensor
clocks and transfer studies. The record also retains failed geometry, clipping,
and nonlinear-readout controls from the range and representation studies.

The multimodal V16 phase-correction family supplied the final-refit recipe.
V18 tested joint measurement, correspondence, geometric readout, predictor
capacity, and encoder adaptation. Some arms improved individual bands, but none
passed every promotion criterion; V18 closed as \texttt{NO\_CANDIDATE}. V19
then tested negative-loss weighting and fixed expert selection.

\subsection{October: final refit and causal runtime}

The final TRAIN40 refit used group-separated auxiliary producers for
out-of-fold inputs and full-TRAIN40 producers for deployment. It followed the
V16 recipe but trained a new model lineage. The submission preserved query-token
identity and complete prediction coverage, and its official score placed
PHASE-TTC first on the public GarlTTC Test leaderboard on 28 October 2026
(Section~\ref{sec:official}).

Runtime experiments covered persistent histories, exact normalization, graph
execution, crop reuse, and multiobject replay. Timing reports separated startup
from measured arrivals and counted valid TTC outputs separately from failed acknowledgements.
The 5\,ms ingestion experiment kept 20\,ms output queries and established
numerical parity for that setting; it did not measure throughput.
